\chapter{Notation}

\section*{Typefaces}

Latin\-/script Interturkic text is set in a serif typeface,
e.g.\ \latntxt{arpa}, \latntxt{fikir}, \latntxt{yılan}.
When embedded within English\-/language prose, Latin\-/script text and
transliterations\---of any foreign language\---are additionally italicised,
e.g.\ Interturkic \foreign{arpa}, Arabic \foreign{kitāb}, Persian \foreign{hēč}.
Glossing abbreviations are set in small caps,
e.g. \Abl, \Spl, \Prs.

Arabic\-/script text\---of any language\---is set in a Naskh\-/like typeface,
e.g.\ \arabtxt{آرْپَا}, \arabtxt{کِتَابْ}, \arabtxt{‌هَیچْ}.
Except for discussion about the Arabic orthography itself, or wherever else it
is deemed necessary, this grammar will use the Latin orthography exclusively.

\section*{Transcriptions}

Phonological transcriptions are rendered with symbols from the International
Phonetic Alphabet~(IPA).
\emph{Phonemic} transcriptions show abstract categories of sound that
fundamentally contribute to an item's meaning,
and are enclosed between forward slashes,
e.g. \phonem{ɑwt͡ʃɯ}, \phonem{ijnɛ}, \phonem{ɔn biɾ}.
By contrast, \emph{phonetic} transcriptions show concrete realisations of sound,
as produced in actual speech,
and are enclosed between square brackets,
e.g.\ \phonet{ɑwd͡ʒɯ}, \phonet{iːne}, \phonet{ɔm biɾ}.

Certain sounds alternate according to morphophonological factors,
and are represented with uppercase letters in both IPA transcriptions and
Latin\-/script text.
Uppercase `A' represents the alternation between \phonem{ɛ} and \phonem{ɑ}.
Uppercase `D' represents the alternation between \phonem{t} and \phonem{d}.
Uppercase `G' represents the alternation between \phonem{k}, \phonem{g}, \phonem
{q}, and \phonem{ɣ}.
Uppercase `I' represents the alternation between \phonem{i}, \phonem{y}, \phonem
{ɯ}, and \phonem{u}.
Uppercase `K' represents the alternation between \phonem{k} and \phonem{q}.
By contrast, lowercase letters represent invariant segments.

\section*{Morphemes}

A dash to the right of an element marks it as a verbal stem,
e.g.\ \foreign{\latntxt{al-}} `take', \foreign{\latntxt{ber-}} `give', \foreign
{\latntxt{kör}} `see'.
A dash to the left of an element marks it as a suffix,
\foreign{\latntxt{\=/GAn}} \intragls{\Ptcp.\Pst}, \foreign{\latntxt{\=/DA}}
\intragls{\Loc}, \foreign{\latntxt{\=/sA}} \intragls{\Cond}.

Parentheses at the beginning of suffix description enclose conditionally
included segments.
An enclosed vowel appears only if preceded by a consonant,
e.g.\ \foreign{\latntxt{üy}} +~\foreign{\latntxt{\=/(I)ñ}} \derivesto~\foreign{%
\latntxt{üy\-/üñ}}.
Conversely, an enclosed consonant appears only if preceded by a vowel,
e.g.\ \foreign{\latntxt{altı}} +~\foreign{\latntxt{\=/(ş)Ar}} \derivesto~%
\foreign{\latntxt{altı\-/şar}}.
Parentheses anywhere else in a suffix description enclose optional segments,
e.g.\ \foreign{\latntxt{öz\-/ü}} +~\foreign{\latntxt{\=/n(I)}} \derivesto~%
\foreign{\latntxt{öz\-/ü\-/nü}} \altwith~\foreign{\latntxt{öz\-/ü\-/n}}.

\section*{Glosses}

Glosses show the morphological breakdown of text fragments.
Morphemes, in both glosses and source text, are separated by dashes,
e.g. \foreign{\latntxt{mez\-/ler\-ge}} \intragls{table\-/\Pl\-/\Dat}.
When a single morpheme is represented by multiple metalanguage elements (i.e.\
words or glossing abbreviations), they are joined together with full stops,
e.g.\ \foreign{\latntxt{\=/(s)I}} \intragls{\Poss.\Third}.

Inline glosses (i.e.\ glosses set inside of prose) are enclosed between angle
brackets,
e.g.\ \intragls{\Gen}, \intragls{\Sg}, \intragls{\Poss.\Third}.
Interlinear glosses are set apart from normal prose and left\-/aligned
vertically per word, e.g.\
\ex
  \intergls
    {Sen Ortaqtürk til\-/i\-/n bil\-/e\-/sen\-/mi?}
    {\Ssg{} common.Turk language\-/\Poss.\Third\-/\Acc{} know\-/\Prs\-/\Ssg\-/%
     \Q}
    {Do you know the Interturkic language?}
\xe

\section*{Other punctuation}

Zero morphemes are represented by \zero.
Single\-/barred arrows express derivation or transformation;
{\derivesto} means `is derived as' and {\derivesfrom} means `is derived from'.
Double\-/barred arrows express copying of foreign elements;
{\copiesto} means `is copied as' and {\copiesfrom} means `is copied as'.
Tildes represent variation between two possibilities with the implication of
both carrying the same meaning,
e.g.\ \phonem{ɔn biɾ} \altwith~\phonet{ɔm biɾ}.